\section*{Chapitre \ref{ch:ll:cpu-microarchitecture} --- Microarchitecture du processeur}

\begin{solution}{exo:ll:cpu-microarchitecture:1}
$600 \times 4.0 = 2\,400$ cycles, soit $2\,400 \times 2.4 = 5\,760$ instructions.
\end{solution}

\begin{solution}{exo:ll:cpu-microarchitecture:2}
$3\,000 / 3.0 = \qty{1\,000}{\nano\second}$, soit $4\,500$ cycles du cœur à 4.5~GHz. Le \omterm{def:ll:cpu-microarchitecture:tsc}{compteur d'horodatage} (TSC) compte le temps.
\end{solution}

\begin{solution}{exo:ll:cpu-microarchitecture:3}
$4 \times 2 = 8$ chaînes (la \cref{prop:ll:cpu-microarchitecture:little})~; trois chaînes achèvent $3/4$ d'opération par
cycle, soit $37.5~\%$ du débit.
\end{solution}

\begin{solution}{exo:ll:cpu-microarchitecture:4}
Avec branchements~: $0.3 \times 20 = 6$ cycles en moyenne~; sans branchement~: 4. La version sans branchement est plus rapide de 2 cycles par
message (la \cref{prop:ll:cpu-microarchitecture:breakeven}~: $m = 0.3 > 4/20$).
\end{solution}

\begin{solution}{exo:ll:cpu-microarchitecture:5}
Sur la mesure retenue, $(2.80 - 0.21)/0.5 = \qty{5.2}{\nano\second}$, soit environ 24 cycles à la fréquence mesurée
4.6~GHz du CPU~2. L'estimation suppose que le prédicteur se trompe sur la moitié des itérations sur des données mélangées (il ne peut pas faire
mieux que le hasard sur des valeurs aléatoires indépendantes) et que rien d'autre ne diffère entre les deux exécutions.
\end{solution}

\begin{solution}{exo:ll:cpu-microarchitecture:6}
Chaque instruction d'une chaîne a besoin du résultat de la précédente, donc au plus une est en vol, quel que soit le nombre d'emplacements dont
dispose le tampon~; la chaîne s'exécute à sa latence. Le tampon masque la latence en recouvrant la chaîne par du travail indépendant (d'autres
chaînes, les chargements des itérations suivantes), et il masque les défauts de cache tant que ce travail tient dans le tampon.
\end{solution}

\begin{solution}{exo:ll:cpu-microarchitecture:7}
\texttt{measured\_blocks.csv}~: les blocs de 1, 2, 4 et 8 tournent à la vitesse du tri (le prédicteur, qui exploite l'historique
récent du branchement, apprend parfaitement les motifs de cette période)~; les blocs de 16 et 64 coûtent environ deux fois plus, une
erreur de prédiction à chaque changement~; les blocs de 1\,024 retrouvent la vitesse du tri, les changements étant rares.
\end{solution}

\begin{solution}{exo:ll:cpu-microarchitecture:8}
Trois défauts. Les chaînes \texttt{add \$1} sont résolues dans le renommeur, environ cinq par cycle sur cette famille de cœurs, si bien que la
référence «~cycle~» est cinq fois trop courte. Un cœur chaud qui exécute un même ordre dix millions de fois garde ses caches et
son prédicteur, parfaitement entraînés, ce que le chemin critique (hot path) ne voit jamais après une période calme. Il faut encadrer la lecture du compteur par des barrières et
en soustraire le coût. Pour l'horloge, utiliser une chaîne d'additions sur registres, faire varier les ordres, et mesurer à froid comme à chaud.
\end{solution}

\begin{solution}{pb:ll:cpu-microarchitecture:1}
Valeurs tirées de la mesure retenue (les \cref{fig:ll:cpu-microarchitecture:branch,fig:ll:cpu-microarchitecture:freq}).
\begin{enumerate}
\item Trié~: avec branchements \qty{0.21}{\nano\second}, sans branchement \qty{0.32}{\nano\second} par élément~; mélangé~: avec branchements
\qty{2.80}{\nano\second}, sans branchement \qty{0.32}{\nano\second}.
\item Le prédicteur a presque toujours raison, et la boucle avec branchements ne fait aucun travail pour la moitié des éléments inférieurs à 128~;
la boucle sans branchement calcule et ajoute un masque pour chaque élément.
\item $(2.80 - 0.21)/0.5 = \qty{5.2}{\nano\second}$.
\item $5.2 \times 4.64 \approx 24$ cycles.
\item Avec branchements~: $0.21 + 5.2\,m$ ns~; sans branchement~: \qty{0.32}{\nano\second}.
\item $m^\ast = 0.11/5.2 \approx 2.1~\%$.
\item Avec branchements, $0.21 + 0.05 = \qty{0.26}{\nano\second}$ contre 0.32~: l'avantage est d'environ \qty{0.06}{\nano\second} par transaction.
\item Sans branchement~: $0.21 + 0.35 \times 5.2 = \qty{2.03}{\nano\second}$ contre environ 0.32~: l'écart est d'environ
\qty{1.7}{\nano\second} par transaction.
\item $50\,000 \times 0.35 \times \qty{5.2}{\nano\second} \approx \qty{91}{\micro\second}$ par seconde.
\item Faible~: \qty{5}{\nano\second} représente environ 0.4~\% de \qty{1.3}{\micro\second}~; le coût tient au partitionnement (clustering),
et non à la moyenne.
\item Les marchés animés sont ceux où le sens de la prochaine transaction est le moins prévisible, et c'est là que la latence
compte le plus.
\item Cela rendrait à nouveau la version avec branchements peu coûteuse~; mais le motif des sens des transactions n'est pas périodique, il y a donc peu
à apprendre.
\item \textbf{Résultat à retenir.} Une erreur de prédiction coûte \qty{5}{\nano\second} sur cet ordinateur portable, soit environ 24 cycles~; la
réécriture sans branchement l'emporte dès que plus d'environ 2~\% des branchements sont mal prédits.
\item Virtualisation, une fréquence inférieure au turbo sous charge, et un indice de référence qui inclut aussi le reste du travail de la boucle~; le
chiffre publié est un minimum sur du matériel nu.
\item Temps écoulé à un rythme constant (environ 3~GHz ici), et non les cycles du cœur, qui suivent son horloge.
\item Pour garder le thread sur un seul cœur avec ses caches et son prédicteur~; sous un hyperviseur, il est épinglé à un CPU virtuel,
qui peut migrer entre les cœurs physiques.
\item Probablement, on transformerait le branchement en déplacement conditionnel ou on vectoriserait la boucle, ce qui supprime l'erreur de prédiction et
l'expérience.
\item Prévisible~: la répartition selon le type de message sur un flux dominé par un seul type. Imprévisible~: savoir si une transaction a eu lieu au
cours acheteur ou au cours vendeur.
\item Avec des compteurs de performance matériel (\texttt{perf stat -e branch-misses}) sur l'hôte de production, ou en comptant
les changements de sens dans les données de marché enregistrées.
\item Quand un branchement dépend de données imprévisibles et que son taux d'erreurs de prédiction dépasse le rapport entre le surcoût
de la réécriture et la pénalité.
\end{enumerate}
\end{solution}

\begin{solution}{iq:ll:cpu-microarchitecture:1}
\omterm{def:ll:cpu-microarchitecture:branch}{Prédiction de branchement}. Sur des données triées, le \texttt{if} va dans le même sens sur de longues séquences et est prédit correctement~; sur des
données aléatoires, il se trompe une fois sur deux, et chaque erreur de prédiction vide le pipeline, ce qui coûte des dizaines de cycles.
\iqlookfor{le vidage du pipeline lors d'une mauvaise prédiction, et le fait que le travail lui-même n'a pas changé.}
\end{solution}

\begin{solution}{iq:ll:cpu-microarchitecture:2}
Le cœur démarre chaque instruction dès que ses entrées sont prêtes, et non dans l'ordre du programme, et valide les résultats dans l'ordre par le
\omterm{def:ll:cpu-microarchitecture:ooo}{tampon de réordonnancement} (reorder buffer). Il est limité par les chaînes de dépendances, par la taille du tampon et du planificateur, par
les \omterm{def:ll:cpu-microarchitecture:branch}{erreurs de prédiction de branchement} (qui annulent le travail) et par la largeur du front-end.
\iqlookfor{ordonnancement piloté par la disponibilité des entrées avec validation dans l'ordre, et les limites.}
\end{solution}

\begin{solution}{iq:ll:cpu-microarchitecture:3}
Un seul accumulateur est une chaîne de dépendances qui s'exécute à la latence d'une addition~; quatre chaînes saturent les additionneurs, dont le
débit vaut la latence multipliée par la largeur (loi de Little). Le compilateur ne peut pas réassocier les additions en virgule flottante, ce qui
modifie le résultat, sauf si on l'y autorise (\texttt{-ffast-math}, ou \texttt{-fassociative-math}).
\iqlookfor{latence et débit, et non-associativité de l'arithmétique en virgule flottante.}
\end{solution}

\begin{solution}{iq:ll:cpu-microarchitecture:4}
L'appeler de nombreuses fois entre des lectures encadrées du \omterm{def:ll:cpu-microarchitecture:tsc}{compteur d'horodatage} (TSC) ou chronométrer un lot, conserver son résultat vivant grâce à une
\omterm{def:ll:cpu-microarchitecture:ubench}{barrière de compilation}, soustraire le surcoût du chronomètre, épingler le thread, faire chauffer le système et donner des percentiles. Pièges~:
l'optimiseur supprimant ou déplaçant le travail, un prédicteur entraîné et des caches chauds, des changements de fréquence, des ticks du TSC pris pour
des cycles.
\iqlookfor{barrières, surcoût, et écart entre conditions à chaud et conditions réalistes.}
\end{solution}

\begin{solution}{iq:ll:cpu-microarchitecture:5}
Non~: le thread frère partage les unités d'exécution, les caches de premier niveau et les prédicteurs, si bien que sa charge ajoute de la latence et de la \omterm{def:ll:where-latency-comes-from:tail}{gigue} (jitter).
Le laisser inactif, ou désactiver le \omterm{def:ll:cpu-microarchitecture:smt}{multithreading simultané} (SMT) sur les cœurs du chemin critique.
\iqlookfor{ressources partagées par cœur et \omterm{def:ll:where-latency-comes-from:tail}{gigue} qu'elles engendrent.}
\end{solution}

\begin{solution}{iq:ll:cpu-microarchitecture:6}
Le cœur sortait d'un \omterm{def:ll:cpu-microarchitecture:dvfs}{état de veille} profond (limiter les C-states et comparer)~; il tourne à basse fréquence avant de monter en charge (figer la
fréquence et comparer)~; les caches, les tampons de traduction et les prédicteurs étaient froids (garder le chemin chaud avec un
trafic synthétique et comparer). Vérifier chacun en chronométrant le premier message avec et sans le correctif.
\iqlookfor{\omterm{def:ll:cpu-microarchitecture:dvfs}{états de veille}, \omterm{def:ll:cpu-microarchitecture:dvfs}{variation dynamique de fréquence}, et état microarchitectural froid, chacun associé à une expérience.}
\end{solution}
